\documentclass[a4paper,11pt]{article}
\usepackage[a4paper,margin=2cm]{geometry}
\usepackage[T1]{fontenc}
\usepackage[utf8]{inputenc}
\usepackage[ngerman,english]{babel}
\usepackage{lmodern,microtype,xcolor,parskip,enumitem}
\usepackage[most]{tcolorbox}
\usepackage{titlesec}
\usepackage[hidelinks]{hyperref}
\definecolor{HeaderBlueDark}{RGB}{70,110,170}
\definecolor{RowBlueLight}{RGB}{235,243,255}
\definecolor{RowBlueDark}{RGB}{220,233,250}
\definecolor{SoftGray}{RGB}{90,90,90}
\setlength{\parindent}{0pt}
\setlist[itemize]{leftmargin=1.4em,itemsep=0.35em,topsep=0.35em}
\titleformat{\section}[block]{\Large\bfseries\color{HeaderBlueDark}}{}{0pt}{}
\newtcolorbox{exbox}{enhanced,breakable,colback=RowBlueLight,colframe=RowBlueDark,boxrule=0.4pt,arc=1.5mm,left=1.2mm,right=1.2mm,top=1mm,bottom=1mm,title={Beispiele \textnormal{(Examples)}},fonttitle=\bfseries,colbacktitle=RowBlueDark,coltitle=black}
\NewDocumentEnvironment{grammarentry}{mm+b}{\begin{tcolorbox}[enhanced,breakable,colback=white,colframe=HeaderBlueDark,title={#1\hfill\normalfont\footnotesize #2},colbacktitle=HeaderBlueDark,coltitle=white,fonttitle=\bfseries]#3\end{tcolorbox}}{}
\newcommand{\GermanText}[1]{\textbf{Deutsch: }#1\par}
\newcommand{\EnglishText}[1]{{\small\color{SoftGray}\textbf{English: }#1\par}}
\newcommand{\ExampleItem}[2]{\item \textbf{#1}\\{\small\color{SoftGray}#2}}
\title{Modalpartikeln / Abtönungspartikeln\\{\large\textnormal{Modal / discourse particles}}}
\date{}
\begin{document}
\maketitle
\tableofcontents
\newpage
\section{Was sind Modalpartikeln? \textnormal{(What are modal particles?)}}
\begin{grammarentry}{Definition}{Modalpartikel / Abtönungspartikel}
\GermanText{Modalpartikeln sind unflektierbare Wörter, die Erwartungen, Bewertungen, Annahmen oder Einstellungen einer sprechenden Person zu einer Äußerung ausdrücken. Sie beeinflussen den Ton, die Beziehung zum Gegenüber und die Gesprächsabsicht. Sie heißen auch Abtönungspartikeln, Einstellungspartikeln oder Diskurspartikeln.}
\EnglishText{Modal particles are invariable words that express the speaker's expectations, assessments, assumptions or attitudes toward an utterance. They affect tone, the relationship with the listener and communicative intent. They are also called discourse particles or attitudinal particles.}
\end{grammarentry}
\GermanText{Meist bleibt die sachliche Kernaussage ohne die Partikel erhalten. Dennoch kann die kommunikative Wirkung stark anders sein: eine Frage klingt neugierig statt neutral, ein Vorschlag freundlicher oder eine Aussage vorwurfsvoll.}
\EnglishText{Usually the basic factual proposition survives without the particle. The communicative effect can nevertheless change significantly: a question can sound curious rather than neutral, a suggestion friendlier or a statement reproachful.}
\begin{exbox}\begin{itemize}
\ExampleItem{Was machst du?}{What are you doing? (neutral question)}
\ExampleItem{Was machst du denn?}{So, what are you doing? (interest or surprise)}
\ExampleItem{Komm mit!}{Come along! (direct command)}
\ExampleItem{Komm doch mit!}{Why don't you come along! (encouragement)}
\ExampleItem{Es ist so.}{That's how it is.}
\ExampleItem{Es ist halt so.}{That's just how it is. (resignation)}
\end{itemize}\end{exbox}
\section{Grammatische Eigenschaften \textnormal{(Grammatical properties)}}
\begin{grammarentry}{Unflektierbar und ohne Rektion}{Morphologie}
\GermanText{Modalpartikeln werden nicht dekliniert oder konjugiert. Sie haben keinen Kasus, kein Genus und keinen Numerus. Sie regieren weder Dativ noch Akkusativ oder Genitiv. Sie verbinden als Modalpartikel auch keine Haupt- oder Nebensätze.}
\EnglishText{Modal particles are neither declined nor conjugated. They have no grammatical case, gender or number. They do not govern dative, accusative or genitive and, when used as modal particles, do not join independent or subordinate clauses.}
\end{grammarentry}
\begin{grammarentry}{Satzbezug statt Satzglied}{Syntax}
\GermanText{Sie beziehen sich normalerweise auf die Äußerung als kommunikative Einheit. Sie sind keine gewöhnlichen Satzglieder: Man kann nicht mit wer?, wo?, wann? oder wie? nach ihnen fragen. Meist lassen sie sich weglassen, ohne den Satz grammatisch unvollständig zu machen. Weglassen kann allerdings Ton und Bedeutung im Gespräch ändern.}
\EnglishText{They normally relate to the utterance as a communicative whole, rather than functioning as ordinary clause constituents. They cannot normally be elicited using who?, where?, when? or how?. They are usually optional for grammatical well-formedness, though removing them may change pragmatic meaning.}
\end{grammarentry}
\begin{grammarentry}{Satzmodussensitiv}{Dependent on sentence mood}
\GermanText{Nicht jede Modalpartikel passt zu jeder Satzart. Manche werden vor allem in Aussagesätzen verwendet, andere in Fragen, Aufforderungen, Wünschen oder Ausrufen. Deshalb reicht es nicht, die deutsche Wortform mit einer einzelnen englischen Bedeutung zu lernen.}
\EnglishText{Not every modal particle fits every sentence mood. Some occur mainly in statements; others in questions, commands, wishes or exclamations. Therefore, learning one fixed English equivalent is insufficient.}
\end{grammarentry}
\section{Position im Satz \textnormal{(Word order and position)}}
\begin{grammarentry}{Mittelfeld}{Standardposition}
\GermanText{Modalpartikeln stehen normalerweise im Mittelfeld. In einem Hauptsatz stehen sie also nach dem finiten Verb und vor einem vorhandenen verbalen Schlussteil. In einem Nebensatz stehen sie vor dem Verbkomplex am Satzende. Sie können im Normalfall nicht allein das Vorfeld einer Aussage besetzen.}
\EnglishText{Modal particles normally appear in the middle field. In a main clause they follow the finite verb and precede any verbal element at the end. In subordinate clauses, they occur before the final verb cluster. They normally cannot independently occupy the prefield of a declarative clause.}
\end{grammarentry}
\begin{exbox}\begin{itemize}
\ExampleItem{Er kommt ja morgen.}{He is coming tomorrow, as you know.}
\ExampleItem{Wann kommt er denn?}{So when is he coming?}
\ExampleItem{Ich weiß nicht, warum er denn nicht kommt.}{I don't know why he isn't coming, then.}
\ExampleItem{Du kannst doch ruhig fragen.}{You can go ahead and ask, you know.}
\ExampleItem{Warum bist du eigentlich hier?}{So why are you here, actually?}
\end{itemize}\end{exbox}
\begin{grammarentry}{Kein Vorfeldtest ohne Kontext}{Different functions of the same word}
\GermanText{Steht ein Wort wie eigentlich, vielleicht oder eben am Satzanfang, kann es in diesem Satz als Adverb gebraucht sein. Man darf aus der Schreibweise allein nicht schließen, dass die Wortform stets Modalpartikel ist. Verkürzte Ausdrücke wie Was denn? sind elliptisch und müssen getrennt von vollständigen Sätzen betrachtet werden.}
\EnglishText{When a word such as eigentlich, vielleicht or eben appears initially, it may be an adverb in that sentence. The written word alone does not determine its grammatical function. Elliptical expressions such as Was denn? need separate analysis.}
\end{grammarentry}
\section{Satzarten und ihre Modalpartikeln \textnormal{(Sentence types and their modal particles)}}
\subsection{Aussagen \textnormal{(Statements)}}
\begin{grammarentry}{Bekanntheit, Widerspruch und Akzeptanz}{Statements}
\GermanText{Ja erinnert an Bekanntes; doch kann einer Erwartung widersprechen; halt und eben drücken oft Akzeptanz oder Unausweichlichkeit aus; schon kann beruhigen; wohl kann eine Vermutung anzeigen.}
\EnglishText{Ja invokes shared knowledge; doch may contradict an expectation; halt and eben often express acceptance or inevitability; schon reassures; wohl may convey supposition.}
\end{grammarentry}
\begin{exbox}\begin{itemize}
\ExampleItem{Das weißt du ja.}{You know that already.}
\ExampleItem{Das stimmt doch!}{But that's true!}
\ExampleItem{Das ist halt so.}{That's just how it is.}
\ExampleItem{Dann musst du eben warten.}{Then you'll just have to wait.}
\ExampleItem{Das wird schon klappen.}{It'll work out.}
\ExampleItem{Das wird wohl stimmen.}{I suppose that's true.}
\end{itemize}\end{exbox}
\subsection{Fragen \textnormal{(Questions)}}
\begin{grammarentry}{Neugier, Überraschung und Zweifel}{Questions}
\GermanText{Denn signalisiert oft Neugier oder Erstaunen; etwa deutet häufig eine unerwartete Möglichkeit an; eigentlich leitet gelegentlich einen Themenwechsel ein; bloß oder nur intensivieren emotionale Fragen; überhaupt kann die grundsätzliche Berechtigung einer Handlung hinterfragen.}
\EnglishText{Denn often conveys interest or surprise; etwa suggests an unexpected possibility; eigentlich can mark a topic shift; bloß or nur intensify emotional questions; überhaupt can challenge the basic reason for an action.}
\end{grammarentry}
\begin{exbox}\begin{itemize}
\ExampleItem{Wo wohnst du denn?}{So where do you live?}
\ExampleItem{Hast du etwa Angst?}{Don't tell me you're afraid?}
\ExampleItem{Wie heißt du eigentlich?}{By the way, what's your name?}
\ExampleItem{Wie soll ich das bloß schaffen?}{How on earth can I manage that?}
\ExampleItem{Warum bist du überhaupt gekommen?}{Why did you come in the first place?}
\end{itemize}\end{exbox}
\subsection{Aufforderungen und Bitten \textnormal{(Commands and requests)}}
\begin{grammarentry}{Einladung, Erlaubnis und Warnung}{Imperatives}
\GermanText{Doch und mal können Bitten freundlicher oder lockerer machen. Ruhig ermutigt oder erlaubt. Bloß und stark betontes ja erzeugen bei Verboten eine dringliche Warnung. Die Wirkung hängt immer auch vom Tonfall ab.}
\EnglishText{Doch and mal can make requests friendlier or more casual. Ruhig encourages or grants permission. Bloß and strongly stressed ja make prohibitions sound urgent. The effect also depends on intonation.}
\end{grammarentry}
\begin{exbox}\begin{itemize}
\ExampleItem{Komm doch mit!}{Why don't you come along!}
\ExampleItem{Schau mal!}{Have a look!}
\ExampleItem{Frag ruhig!}{Go ahead and ask!}
\ExampleItem{Mach bloß keinen Unsinn!}{Whatever you do, don't do anything foolish!}
\ExampleItem{Fass das ja nicht an!}{Don't you dare touch that!}
\end{itemize}\end{exbox}
\subsection{Wünsche und Ausrufe \textnormal{(Wishes and exclamations)}}
\begin{grammarentry}{Sehnsucht und starke Bewertung}{Wishes and exclamations}
\GermanText{Doch und nur können einen unerfüllten Wunsch betonen. Aber, ja und vielleicht können überraschte oder emotionale Ausrufe verstärken. Vielleicht bedeutet in diesem Ausrufgebrauch nicht möglicherweise.}
\EnglishText{Doch and nur can intensify an unfulfilled wish. Aber, ja and vielleicht can strengthen surprised or emotional exclamations. In this exclamatory use, vielleicht does not mean perhaps.}
\end{grammarentry}
\begin{exbox}\begin{itemize}
\ExampleItem{Wenn ich doch nur mehr Zeit hätte!}{If only I had more time!}
\ExampleItem{Das ist aber schön!}{Oh, that's lovely!}
\ExampleItem{Du bist ja schon hier!}{Oh, you're already here!}
\ExampleItem{Das war vielleicht ein Chaos!}{What a mess that was!}
\end{itemize}\end{exbox}
\section{Betonung und Sprechweise \textnormal{(Stress and intonation)}}
\begin{grammarentry}{Unbetont ist häufig, aber nicht immer}{Stressed particles}
\GermanText{Viele Modalpartikeln sind unbetont. Einige haben besondere betonte Verwendungen, beispielsweise ja in Pass JA auf! oder bloß in Mach BLOSS keinen Fehler! (Betonungsmarkierung durch Großschreibung). Betonung kann die kommunikative Wirkung deutlich verschärfen.}
\EnglishText{Many modal particles are unstressed. Some have distinctive stressed uses, such as ja in Pass JA auf! or bloß in Mach BLOSS keinen Fehler! (capitals here mark emphasis). Stress can make the utterance sound much more forceful.}
\end{grammarentry}
\GermanText{Entscheidend sind nicht nur die Wörter, sondern Gesprächssituation, Satzart und Intonation. Beispielsweise kann Was machst du denn? freundlich, erstaunt oder vorwurfsvoll klingen.}
\EnglishText{What matters is not just the words, but the situation, sentence mood and intonation. For example, Was machst du denn? can sound friendly, surprised or reproachful.}
\section{Mehrere Modalpartikeln in einem Satz \textnormal{(Combining modal particles)}}
\begin{grammarentry}{Kombination und Reihenfolge}{Ordering restrictions}
\GermanText{Mehrere Modalpartikeln dürfen zusammenstehen. Die Reihenfolge ist eingeschränkt und hängt besonders von der Satzart ab. In Fragen ist denn vor etwa oder bloß typisch; in Aufforderungen steht doch oft vor mal oder ruhig. Das ist keine Regel, mit der jede beliebige Kombination erlaubt wäre.}
\EnglishText{Several modal particles may occur together. Their order is restricted and strongly depends on sentence mood. In questions, denn commonly precedes etwa or bloß; in commands, doch often precedes mal or ruhig. This does not license arbitrary combinations.}
\end{grammarentry}
\begin{exbox}\begin{itemize}
\ExampleItem{Was ist denn bloß passiert?}{What on earth happened?}
\ExampleItem{Hast du denn etwa alles vergessen?}{Have you seriously forgotten everything?}
\ExampleItem{Komm doch einfach mal vorbei!}{Why don't you just drop by sometime!}
\ExampleItem{Frag doch ruhig mal nach!}{Do feel free to ask!}
\ExampleItem{Du weißt ja doch, was ich meine.}{You do know what I mean, after all.}
\end{itemize}\end{exbox}
\section{Die Verbindungen denn nicht und etwa nicht \textnormal{(The sequences denn nicht and etwa nicht)}}
\begin{grammarentry}{denn nicht}{Modalpartikel + Negationspartikel}
\GermanText{In Warum kommst du denn nicht? steht denn als Modalpartikel und nicht als Negationspartikel. Denn verändert die Gesprächswirkung; nicht verneint den Satzinhalt. Beide Wörter können durch andere Satzglieder getrennt werden und bilden keine einzelne Modalpartikel.}
\EnglishText{In Warum kommst du denn nicht?, denn is a modal particle and nicht a negation particle. Denn modifies conversational force, while nicht negates the proposition. They may be separated by other constituents and do not form one compound modal particle.}
\end{grammarentry}
\begin{exbox}\begin{itemize}
\ExampleItem{Warum kommst du denn nicht?}{Why aren't you coming, then?}
\ExampleItem{Warum hast du denn gestern nicht angerufen?}{Why didn't you call yesterday, then?}
\ExampleItem{Warum kommst du nicht?}{Why aren't you coming? (more neutral without denn)}
\end{itemize}\end{exbox}
\begin{grammarentry}{etwa nicht}{Unerwartete Verneinung}
\GermanText{In Kommst du etwa nicht? ist etwa die Modalpartikel und nicht die Negationspartikel. Der Sprecher hat häufig erwartet, dass die Person kommt. Durch etwa wird die gegenteilige Möglichkeit als überraschend oder unerwünscht dargestellt.}
\EnglishText{In Kommst du etwa nicht?, etwa is the modal particle and nicht is the negation particle. The speaker often expected the person to come. Etwa makes the contrary possibility sound unexpected or unwelcome.}
\end{grammarentry}
\begin{exbox}\begin{itemize}
\ExampleItem{Kommst du etwa nicht?}{Don't tell me you're not coming?}
\ExampleItem{Hast du etwa nicht gelernt?}{Don't tell me you didn't study?}
\ExampleItem{Kommst du denn etwa nicht?}{So, are you seriously not coming?}
\end{itemize}\end{exbox}
\begin{grammarentry}{Nicht als Sonderfall}{Rare non-negating nicht}
\GermanText{Nicht ist normalerweise eine Negationspartikel. Nur in bestimmten Ausrufen und rhetorischen Fragen kann nicht ohne gewöhnliche negative Bedeutung auftreten; manche Grammatiken zählen solche Fälle zu den Abtönungspartikeln. Das gilt nicht automatisch für denn nicht oder etwa nicht.}
\EnglishText{Nicht is normally a negation particle. Only in certain exclamations and rhetorical questions can it appear without normal negative force; some grammars classify those uses as modal. This does not automatically apply to denn nicht or etwa nicht.}
\end{grammarentry}
\begin{exbox}\begin{itemize}
\ExampleItem{Was sie nicht alles weiß!}{It's amazing how much she knows!}
\ExampleItem{Was wir nicht alles erlebt haben!}{Just think of everything we have experienced!}
\ExampleItem{Ich komme nicht.}{I'm not coming. (ordinary negation)}
\end{itemize}\end{exbox}
\section{Gleiche Wortform, andere Wortart \textnormal{(Same spelling, different word class)}}
\begin{grammarentry}{denn}{Konjunktion versus Modalpartikel}
\GermanText{Denn kann eine koordinierende Kausalkonjunktion mit der Bedeutung weil sein: Ich bleibe zu Hause, denn ich bin krank. Als Modalpartikel in Was machst du denn? verbindet es keine zwei Hauptsätze.}
\EnglishText{Denn can be a coordinating conjunction meaning because: I am staying home because I am ill. As a particle in Was machst du denn?, it does not join clauses.}
\end{grammarentry}
\begin{grammarentry}{etwa, nur, auch}{Grad- und Fokuspartikeln}
\GermanText{Etwa zwanzig Euro enthält eine ungefähre Zahl; nur zwei Euro begrenzt die Zahl; auch Maria bedeutet Maria ebenfalls. In diesen Fällen liegt keine typische Abtönungsfunktion vor.}
\EnglishText{Etwa zwanzig Euro approximates a number; nur zwei Euro limits it; auch Maria means Maria too. These are not typical modal-particle uses.}
\end{grammarentry}
\begin{grammarentry}{ja, doch, schon, vielleicht}{Antwort, Zeit und Möglichkeit}
\GermanText{Ja! ist eine Antwortpartikel; Doch! kann eine verneinte Frage positiv beantworten; schon in Er ist schon da ist zeitlich; vielleicht in Vielleicht kommt er drückt Möglichkeit aus. Eine Wortform kann also mehrere grammatische Aufgaben haben.}
\EnglishText{Ja! is an answer particle; Doch! can give an affirmative reply to a negative question; schon in Er ist schon da is temporal; vielleicht in Vielleicht kommt er expresses possibility. One word form can thus serve different grammatical roles.}
\end{grammarentry}
\begin{grammarentry}{Modalpartikel versus Modalverb}{No conjugation or infinitive requirement}
\GermanText{Modalverben wie können, müssen, dürfen, sollen, wollen und mögen sind Verben und werden konjugiert. Modalpartikeln wie denn, doch und ja sind dagegen unveränderliche Partikeln. Modalpartikeln verlangen keinen Infinitiv und bilden kein Prädikat.}
\EnglishText{Modal verbs such as können, müssen, dürfen, sollen, wollen and mögen are conjugated verbs. Modal particles such as denn, doch and ja are invariable. They require no infinitive and do not form the predicate.}
\end{grammarentry}
\section{Regionale Verwendungen \textnormal{(Regional usages)}}
\begin{grammarentry}{Österreich, Süd- und Norddeutschland}{Regional particles}
\GermanText{In Österreich hört man häufig eh (Das weißt du eh). In süddeutschen Dialekten kommt fei vor (Das ist fei wichtig). Im norddeutschen Sprachraum gibt es umgangssprachliches man als Abtönungspartikel (Lass man gut sein). Diese regionalen Formen sind nicht für die neutrale schriftliche ÖSD-B1-Produktion erforderlich.}
\EnglishText{In Austria, eh is common (You know that anyway). Southern German dialects use fei (Mind you, that's important). Northern colloquial German may use man as a particle (Come on, let it be). These regional forms are not required in neutral written ÖSD B1 production.}
\end{grammarentry}
\section{Häufige Missverständnisse \textnormal{(Common misunderstandings)}}
\begin{exbox}\begin{itemize}
\ExampleItem{Falsch als Modalpartikel: Denn, was machst du?}{Incorrect if denn is intended as an ordinary fronted modal particle.}
\ExampleItem{Richtig: Was machst du denn?}{Correct: So, what are you doing?}
\ExampleItem{Etwa zwanzig Minuten.}{About twenty minutes. (approximation, not modal particle)}
\ExampleItem{Hast du etwa Angst?}{Don't tell me you're afraid? (modal particle)}
\ExampleItem{Er ist schon da.}{He's already here. (temporal)}
\ExampleItem{Das wird schon klappen.}{It'll work out. (reassurance)}
\ExampleItem{Ich habe nur drei Euro.}{I only have three euros. (focus particle)}
\ExampleItem{Wie soll ich das nur machen?}{How on earth can I do that? (modal particle)}
\end{itemize}\end{exbox}
\GermanText{Eine ungewöhnliche Modalpartikel ist nicht unbedingt ein Grammatikfehler. Sie kann jedoch ungewollt Vorwurf, Ironie oder Ungeduld signalisieren. In Prüfungen ist ein klarer neutraler Satz besser als eine unsicher verwendete Partikel.}
\EnglishText{An unusual particle does not automatically make a sentence ungrammatical. It may, however, unintentionally suggest blame, irony or impatience. In an exam, a clear neutral sentence is preferable to a poorly controlled particle.}
\section{Verzeichnis der Einzeldateien \textnormal{(Index of individual particle files)}}
\begin{itemize}
\item \texttt{part/aber.tex} -- Surprise
\item \texttt{part/auch.tex} -- Rhetorical questions
\item \texttt{part/bloss.tex} -- Warnings, wishes, emotional questions (bloß)
\item \texttt{part/denn.tex} -- Questions and surprise
\item \texttt{part/doch.tex} -- Contradiction, encouragement
\item \texttt{part/eben.tex} -- Inevitability
\item \texttt{part/einfach.tex} -- Simple solution
\item \texttt{part/eigentlich.tex} -- Topic shifting
\item \texttt{part/etwa.tex} -- Unexpected possibility
\item \texttt{part/halt.tex} -- Acceptance
\item \texttt{part/ja.tex} -- Shared knowledge, surprise, warning
\item \texttt{part/mal.tex} -- Casual requests
\item \texttt{part/nicht.tex} -- Rare non-negative exclamation
\item \texttt{part/nur.tex} -- Wishes and helplessness
\item \texttt{part/ruhig.tex} -- Permission and encouragement
\item \texttt{part/schon.tex} -- Reassurance
\item \texttt{part/ueberhaupt.tex} -- Fundamental challenge (überhaupt)
\item \texttt{part/vielleicht.tex} -- Emotional exclamation
\item \texttt{part/wohl.tex} -- Supposition
\item \texttt{part/eh.tex} -- Austrian regional usage
\item \texttt{part/fei.tex} -- Southern regional usage
\item \texttt{part/man.tex} -- Northern regional usage
\end{itemize}
\GermanText{Alle Dateien sind eigenständige LaTeX-Erklärungen nach der deutschen Erklärungsvorlage. Für Wörter mit mehreren Wortarten wird jeweils ausdrücklich zwischen der Modalpartikel und einer anderen Funktion unterschieden.}
\EnglishText{Each file is a self-contained LaTeX explanation based on the German explanation template. Where a word has multiple grammatical functions, its modal use is explicitly distinguished from other uses.}
\section{Kurze Übungen \textnormal{(Short exercises)}}
\GermanText{Bestimmen Sie, ob das fett gedruckte Wort hier eine Modalpartikel, eine Konjunktion, eine Negationspartikel, eine Fokuspartikel oder ein anderes Adverb ist.}
\EnglishText{Decide whether the word in bold is a modal particle, a conjunction, a negation particle, a focus particle or another adverb.}
\begin{exbox}\begin{enumerate}
\item Warum kommst du \textbf{denn} \textbf{nicht}?
\item Ich bleibe hier, \textbf{denn} ich bin krank.
\item Hast du \textbf{etwa} die Prüfung vergessen?
\item Das kostet \textbf{etwa} hundert Euro.
\item Wie soll ich das \textbf{nur} schaffen?
\item Ich habe \textbf{nur} fünf Euro.
\item Das ist \textbf{ja} interessant!
\item \textbf{Ja}, ich komme.
\item Er ist \textbf{schon} hier.
\item Das wird \textbf{schon} klappen.
\end{enumerate}\end{exbox}
\subsection{Lösungen \textnormal{(Answers)}}
\GermanText{1. denn: Modalpartikel; nicht: Negationspartikel. 2. Konjunktion. 3. Modalpartikel. 4. ungefähre Mengenangabe. 5. Modalpartikel. 6. Fokuspartikel. 7. Modalpartikel. 8. Antwortpartikel. 9. temporales Adverb. 10. Modalpartikel.}
\EnglishText{1. denn: modal particle; nicht: negation particle. 2. Conjunction. 3. Modal particle. 4. Approximate quantity modifier. 5. Modal particle. 6. Focus particle. 7. Modal particle. 8. Answer particle. 9. Temporal adverb. 10. Modal particle.}
\section{Fachquellen \textnormal{(Reference grammar)}}
\GermanText{Die genaue Größe der Klasse der Abtönungspartikeln unterscheidet sich je nach Grammatikmodell; bei auch, einfach, eigentlich, wohl, überhaupt und nicht ist die Grenzziehung besonders wichtig. Für eine vertiefte Analyse helfen die folgenden Quellen.}
\EnglishText{The precise membership of the modal-particle class varies among grammar frameworks; borderline cases especially include auch, einfach, eigentlich, wohl, überhaupt and nicht. The following sources provide further explanation.}
\begin{itemize}
\item IDS Mannheim, \emph{Abtönungspartikel}: \url{https://grammis.ids-mannheim.de/terminologie/2}.
\item IDS Mannheim, \emph{Abtönungspartikeln und Wortstellung}: \url{https://grammis.ids-mannheim.de/systematische-grammatik/769}.
\item Duden, \emph{Modalpartikeln}: \url{https://www.duden.de/sprachwissen/sprachratgeber/Modalpartikeln}.
\item LEO-Grammatik, \emph{Die Abtönungspartikeln}: \url{https://dict.leo.org/grammatik/deutsch/Wort/Adverb/Partikel/Modalpartikel.xml?lang=de}.
\end{itemize}
\end{document}
